\documentclass[12pt,a4paper]{article}
\usepackage[utf8]{inputenc}
\usepackage[T1]{fontenc}
\usepackage[brazil]{babel}
\usepackage{indentfirst}
\usepackage[margin=2.5cm]{geometry}

\title{Resenha Cr\'itica: Artigo 13 ---\\
Object Constraint Language: A Definitive Guide}
\author{Gustavo Pessoa Firmino Duarte}
\date{23 de Abril de 2026}

\begin{document}
\maketitle

\section{Introdu\c{c}\~ao}

A \textit{Object Constraint Language} (OCL) \'e uma linguagem formal de especifica\c{c}\~ao
padronizada pelo \textit{Object Management Group} (OMG) como complemento ao UML
(\textit{Unified Modeling Language}). O guia definitivo da OCL apresenta a linguagem
como uma forma precisa e n\~ao-ambigu\'ia de expressar restri\c{c}\~oes sobre
modelos UML --- restri\c{c}\~oes que n\~ao podem ser expressas adequadamente
atrav\'es de diagramas gr\'aficos. Por meio de express\~oes OCL \'e poss\'ivel
especificar invariantes de classe, pr\'e-condi\c{c}\~oes e p\'os-condi\c{c}\~oes
de opera\c{c}\~oes (conectando-se diretamente ao \textit{Design by Contract} de
Meyer), regras de deriva\c{c}\~ao de atributos e guardas de transi\c{c}\~oes de
estados. A OCL \'e uma linguagem puramente declarativa e sem efeitos colaterais:
avaliar uma express\~ao OCL nunca modifica o estado do sistema.

\section{Sintaxe, Tipos e Principais Construtores}

A OCL possui um sistema de tipos rico que inclui tipos b\'asicos (Boolean, Integer,
Real, String), tipos de cole\c{c}\~ao (Set, OrderedSet, Bag, Sequence) e tipos
definidos pelo modelo UML (as pr\'oprias classes do diagrama de classes). As
express\~oes s\~ao escritas sobre o contexto de um elemento do modelo --- por
exemplo, uma invariante \'e sempre associada a um contexto de classe.

Um exemplo cl\'assico de invariante OCL \'e a restri\c{c}\~ao de que a idade de
uma pessoa deve ser positiva: \texttt{context Pessoa inv: self.idade > 0}. Uma
pr\'e-condi\c{c}\~ao em uma opera\c{c}\~ao de saque banc\'ario poderia ser:
\texttt{context Conta::sacar(valor : Real) pre: valor > 0 and saldo >= valor}.

As opera\c{c}\~oes sobre cole\c{c}\~oes s\~ao especialmente poderosas: \texttt{select},
\texttt{reject}, \texttt{collect}, \texttt{forAll} e \texttt{exists} permitem
expressar restri\c{c}\~oes complexas sobre conjuntos de objetos relacionados de
forma concisa e leg\'ivel. Por exemplo, ``todos os funcion\'arios de um departamento
devem ter sal\'ario maior que o sal\'ario m\'inimo'' se expressa naturalmente em
OCL sem ambiguidade.

A OCL tamb\'em suporta navega\c{c}\~ao atrav\'es de associa\c{c}\~oes UML,
permitindo especificar restri\c{c}\~oes que envolvem m\'ultiplas classes relacionadas.
Esse recurso \'e essencial para modelar regras de neg\'ocio que naturalmente
atravessam fronteiras de objetos.

\section{Conex\~ao com a Pr\'atica Profissional}

A OCL representa o extremo formal do espectro de especifica\c{c}\~ao --- o oposto
do c\'odigo informal que \'e a norma nos projetos em que participei. No meu MVP
de e-commerce, as regras de neg\'ocio como ``um pedido n\~ao pode ser submetido
com carrinho vazio'' ou ``o desconto n\~ao pode exceder 100\% do pre\c{c}o'' foram
implementadas diretamente no c\'odigo de servi\c{c}o, sem especifica\c{c}\~ao
formal pr\'evia. A OCL seria a linguagem adequada para capturar essas regras de
forma inequ\'ivoca antes da implementa\c{c}\~ao.

Conectando com minha leitura do \textit{Design by Contract} de Meyer, percebo
que a OCL \'e essencialmente uma implementa\c{c}\~ao formal do DbC no universo
de modelagem UML: pr\'e-condi\c{c}\~oes, p\'os-condi\c{c}\~oes e invariantes
t\^em mapeamento direto. A diferen\c{c}a \'e que a OCL opera sobre modelos,
enquanto o DbC de Meyer opera diretamente no c\'odigo Eiffel. Em ambientes onde
a modelagem UML \'e parte central do processo de desenvolvimento (como em alguns
contextos de engenharia de sistemas ou conformidade regulat\'oria), a OCL \'e
insubstitu\'ivel como ferramenta de especifica\c{c}\~ao.

\section{Conclus\~ao}

A OCL representa um marco na especifica\c{c}\~ao formal de software orientado a
objetos. Sua integra\c{c}\~ao com UML permite que modelos de an\'alise sejam
enriquecidos com restri\c{c}\~oes precisas, reduzindo a ambiguidade que frequentemente
leva a implementa\c{c}\~oes incorretas. Embora seu uso na ind\'ustria seja limitado
--- a maioria das equipes prefere express\~oes impl\'icitas no c\'odigo ou em
linguagem natural --- a OCL permanece relevante em dom\'inios de alta criticidade,
como sistemas m\'edicos, financeiros e aeroespaciais, onde a especifica\c{c}\~ao
formal \'e uma exig\^encia de conformidade. O estudo da OCL tamb\'em tem valor
pedag\'ogico: for\c{c}a o desenvolvedor a pensar com rigor l\'ogico sobre as
propriedades que o sistema deve satisfazer, uma habilidade transfer\'ivel para
qualquer contexto de desenvolvimento.

\end{document}
